\documentclass[11pt,twocolumn]{article}

% ── Page layout ────────────────────────────────────────────────────────────────
\usepackage[margin=0.58in,top=0.52in,bottom=0.55in,columnsep=0.22in]{geometry}
\usepackage[T1]{fontenc}
\usepackage[utf8]{inputenc}
\usepackage{lmodern}
\usepackage{microtype}
\usepackage{amsmath,amssymb}
\usepackage{graphicx}
\usepackage{booktabs}
\usepackage{caption}
\usepackage{subcaption}
\usepackage{float}
\usepackage{hyperref}
\usepackage{xcolor}
\usepackage{titlesec}
\usepackage{enumitem}
\usepackage[numbers,sort&compress]{natbib}

\hypersetup{colorlinks=true,linkcolor=blue,citecolor=blue,urlcolor=blue}
\setlist{nosep,leftmargin=*}
\captionsetup{font=small,labelfont=bf,skip=2pt}
\captionsetup[subfigure]{font=small,labelfont=bf,skip=1pt}

% Compact but still readable 11pt body text.
\setlength{\parindent}{0pt}
\setlength{\parskip}{1.5pt}
\setlength{\textfloatsep}{5pt plus 1pt minus 1pt}
\setlength{\floatsep}{4pt plus 1pt minus 1pt}
\setlength{\intextsep}{4pt plus 1pt minus 1pt}
\setlength{\abovedisplayskip}{3pt}
\setlength{\belowdisplayskip}{3pt}
\titlespacing*{\section}{0pt}{5pt}{2pt}
\titlespacing*{\subsection}{0pt}{3pt}{1pt}
\titleformat{\section}{\normalfont\bfseries\normalsize}{\thesection.}{0.45em}{}
\titleformat{\subsection}{\normalfont\bfseries\small}{\thesubsection.}{0.45em}{}

% Change this if your figures live somewhere else.
\newcommand{\figdir}{figures}

% Allows the report to compile even if a figure is missing locally.
\newcommand{\maybeincludegraphics}[2][]{%
  \IfFileExists{#2}{\includegraphics[#1]{#2}}{%
    \fbox{\parbox[c][1.08in][c]{0.92\linewidth}{\centering\small Missing figure:\\[-1pt]\texttt{\detokenize{#2}}}}%
  }%
}

\begin{document}

\twocolumn[{%
\centering
{\Large\bfseries Reimplementing LoRA: Low-Rank Adaptation of Large Language Models\par}
\vspace{2pt}
{\normalsize Rohan Shankar\textsuperscript{1} \quad Harshaan Chugh\textsuperscript{1}\par}
{\small \textsuperscript{1}Cornell University, CS 4782 --- Spring 2026 \quad
\texttt{\{rs2656,hsc53\}@cornell.edu}\par}
\vspace{3pt}
\rule{\textwidth}{0.4pt}
\vspace{3pt}
}]

\noindent\textbf{Abstract.}
We reimplemented LoRA~\cite{hu2022lora} from scratch and evaluated it on SST-2 and MRPC with BERT-base and on WikiText-2 with GPT-2 Small. LoRA achieved \textbf{93.1\%} validation accuracy on SST-2 and \textbf{89.9\%} validation F1 on MRPC while training only 0.27\% of parameters. On WikiText-2, LoRA reached 24.27 perplexity versus 23.60 for full fine-tuning while using 99.5\% fewer trainable parameters. Overall, our results support LoRA's qualitative claim: low-rank task updates can preserve strong performance with dramatically fewer trainable weights.

\section{Introduction}
Large pretrained language models are expensive to adapt because full fine-tuning updates every parameter and usually requires storing a separate full model copy per task. We reproduce \emph{LoRA: Low-Rank Adaptation of Large Language Models} by Hu, Shen, Wallis, Allen-Zhu, Li, Wang, Wang, and Chen~\cite{hu2022lora}. The paper observes that downstream weight updates often have low intrinsic rank and proposes freezing pretrained weights while learning a low-rank additive update for selected matrices.

Our project reimplements LoRA directly rather than using PEFT or adapter libraries. We test whether the same efficiency-quality tradeoff appears in a smaller reproducible setting: BERT-base~\cite{devlin2019bert} on GLUE tasks SST-2 and MRPC~\cite{wang2019glue,socher2013sst,dolan2005mrpc}, and GPT-2 Small~\cite{radford2019language} on WikiText-2~\cite{merity2017wikitext}.

\section{Chosen Result}
We aimed to reproduce the paper's central empirical result: LoRA can match or approach full fine-tuning performance while training only a small number of low-rank parameters. This result is tied to the paper's low-rank update equation, $\Delta W=BA$, and its benchmark comparisons showing competitive adaptation with far fewer trainable weights~\cite{hu2022lora}.

We chose this result because it tests LoRA's main scientific claim rather than only a surface-level implementation detail. If task adaptation is well approximated by a low-rank update, then LoRA should remain competitive with full fine-tuning, performance should improve with rank until saturation, and the trainable-parameter budget should be much smaller.

\section{Methodology}
\subsection{LoRA Layer Design}
For a frozen pretrained weight matrix $W_0\in\mathbb{R}^{d_{out}\times d_{in}}$, LoRA learns two small matrices $A\in\mathbb{R}^{r\times d_{in}}$ and $B\in\mathbb{R}^{d_{out}\times r}$:
\begin{equation}
  h = W_0x + \frac{\alpha}{r}BAx.
  \label{eq:lora}
\end{equation}
Here $r$ is the LoRA rank and $\alpha/r$ scales the update. We initialize $A$ with Kaiming uniform initialization and $B$ to zeros so that $\Delta W=0$ at the start of training. This preserves the base model's initial behavior before task-specific learning begins.

We implemented \texttt{LoRALinear} for standard \texttt{nn.Linear} modules and \texttt{LoRAConv1D} for GPT-2's \texttt{Conv1D} attention projection. Both wrappers compute the frozen base path and the low-rank residual path in parallel, then add the scaled residual.

\subsection{Injection, Freezing, and Models}
Our \texttt{apply\_lora\_to\_model()} routine traverses the model graph and replaces modules whose names match target suffixes. For BERT-base, LoRA is inserted into the \texttt{query} and \texttt{value} projections in all transformer layers. For GPT-2 Small, LoRA is inserted into \texttt{c\_attn}, GPT-2's combined attention projection. After replacement, pretrained parameters are frozen; only \texttt{lora\_A}, \texttt{lora\_B}, and, for classification, the task head are trainable.

\subsection{Datasets, Metrics, and Training Setup}
For classification, we evaluate SST-2 using validation accuracy and MRPC using validation F1. For language modeling, we evaluate WikiText-2 using validation perplexity. We vary $r\in\{1,4,8,16\}$ to test the low-rank capacity hypothesis.

For BERT-base experiments, LoRA is applied to \texttt{query,value}. For GPT-2 Small experiments, LoRA is applied to \texttt{c\_attn}. Our scripts and notebooks support these configurations and produce the reported tables and figures. Because our experiments used limited Colab compute, this is a smaller-scale reproduction: the goal is to reproduce the qualitative trend---competitive performance with far fewer trainable parameters---rather than exactly match every original paper number.

\section{Results \& Analysis}
Table~\ref{tab:results} summarizes the best validation metrics. On SST-2 and MRPC, LoRA outperformed our full fine-tuning baselines: $r=8$ achieved 93.1\% accuracy on SST-2 and 89.9\% F1 on MRPC while training only 0.27\% of parameters. This supports the main LoRA claim that constrained low-rank adaptation can be highly parameter-efficient.

\begin{table}[H]
\centering
\caption{Best validation metrics. Parentheses show percent trainable parameters.}
\label{tab:results}
\setlength{\tabcolsep}{2.5pt}
\renewcommand{\arraystretch}{0.98}
{\small
\begin{tabular}{@{}llrl@{}}
\toprule
\textbf{Task} & \textbf{Method} & \textbf{Params} & \textbf{Metric} \\
\midrule
SST-2 & Full FT & 109.5M & 85.3\% acc \\
      & LoRA $r=1$  & 38K (0.04\%) & 92.0\% acc \\
      & LoRA $r=4$  & 149K (0.14\%) & 92.7\% acc \\
      & LoRA $r=8$  & 296K (0.27\%) & \textbf{93.1\% acc} \\
      & LoRA $r=16$ & 591K (0.54\%) & 92.7\% acc \\
\midrule
MRPC  & Full FT & 109.5M & 82.6\% F1 \\
      & LoRA $r=1$  & 38K (0.04\%) & 88.8\% F1 \\
      & LoRA $r=4$  & 149K (0.14\%) & 88.6\% F1 \\
      & LoRA $r=8$  & 296K (0.27\%) & \textbf{89.9\% F1} \\
      & LoRA $r=16$ & 591K (0.54\%) & 89.6\% F1 \\
\midrule
Wiki-2 & Full FT & 124.4M & \textbf{23.60 ppl} \\
       & LoRA $r=1$  & 37K (0.03\%) & 24.99 ppl \\
       & LoRA $r=4$  & 147K (0.12\%) & 24.43 ppl \\
       & LoRA $r=8$  & 295K (0.24\%) & 24.34 ppl \\
       & LoRA $r=16$ & 590K (0.47\%) & 24.27 ppl \\
\bottomrule
\end{tabular}}
\end{table}

The main discrepancy is that our full fine-tuning baselines underperformed LoRA on the GLUE tasks. We attribute this to overfitting and optimization sensitivity: full fine-tuning updates 109.5M parameters on relatively small classification datasets, while LoRA restricts learning to a lower-dimensional subspace and acts as an implicit regularizer.

On WikiText-2, full fine-tuning remained slightly better: 23.60 perplexity versus 24.27 for LoRA at $r=16$. This is reasonable because language modeling may benefit from updating more of the model, and our GPT-2 setup only targeted \texttt{c\_attn}. Still, the gap was small relative to the parameter savings.

The rank ablation supports the low-rank hypothesis. Classification performance improves from $r=1$ to $r=8$ and then saturates or slightly declines at $r=16$. For WikiText-2, perplexity improves steadily through $r=16$. This suggests that useful rank depends on task type and scale: small classification tasks need few adaptation degrees of freedom, while language modeling benefits from more capacity. LoRA also reduced peak GPU memory by up to 58\% on SST-2, reinforcing its practical value.

\section{Reflections}
This project taught us that parameter-efficient fine-tuning depends on both the mathematical idea and the implementation details. The most important choices were which modules to replace, whether base weights were truly frozen, how $A$ and $B$ were initialized, and whether the task head remained trainable.

Our results reproduce LoRA's key trend: low-rank updates can achieve strong downstream performance with orders of magnitude fewer trainable parameters. Given more time, we would run multiple seeds, tune full fine-tuning baselines more carefully, test additional insertion targets such as key/output/MLP projections, and add LoRA weight merging to measure inference latency directly. We would also explore rank-adaptive methods that learn a per-layer rank automatically rather than fixing $r$ as a global hyperparameter.

{
\begin{thebibliography}{10}
\setlength{\itemsep}{-1pt}
\setlength{\parskip}{0pt}

\bibitem[1]{hu2022lora}
E.~J. Hu, Y. Shen, P. Wallis, Z. Allen-Zhu, Y. Li, S. Wang, L. Wang, and W. Chen. LoRA: Low-rank adaptation of large language models. \emph{ICLR}, 2022.

\bibitem[2]{devlin2019bert}
J. Devlin, M.-W. Chang, K. Lee, and K. Toutanova. BERT: Pre-training of deep bidirectional transformers for language understanding. \emph{NAACL-HLT}, 2019.

\bibitem[3]{radford2019language}
A. Radford, J. Wu, R. Child, D. Luan, D. Amodei, and I. Sutskever. Language models are unsupervised multitask learners. OpenAI technical report, 2019.

\bibitem[4]{wang2019glue}
A. Wang, A. Singh, J. Michael, F. Hill, O. Levy, and S. Bowman. GLUE: A multi-task benchmark and analysis platform for natural language understanding. \emph{ICLR}, 2019.

\bibitem[5]{socher2013sst}
R. Socher, A. Perelygin, J. Wu, J. Chuang, C. Manning, A. Ng, and C. Potts. Recursive deep models for semantic compositionality over a sentiment treebank. \emph{EMNLP}, 2013.

\bibitem[6]{dolan2005mrpc}
W. Dolan and C. Brockett. Automatically constructing a corpus of sentential paraphrases. \emph{IWP}, 2005.

\bibitem[7]{merity2017wikitext}
S. Merity, C. Xiong, J. Bradbury, and R. Socher. Pointer sentinel mixture models. \emph{ICLR}, 2017.

\bibitem[8]{paszke2019pytorch}
A. Paszke et al. PyTorch: An imperative style, high-performance deep learning library. \emph{NeurIPS}, 2019.

\bibitem[9]{wolf2020transformers}
T. Wolf et al. Transformers: State-of-the-art natural language processing. \emph{EMNLP: System Demonstrations}, 2020.

\bibitem[10]{lhoest2021datasets}
Q. Lhoest et al. Datasets: A community library for natural language processing. \emph{EMNLP: System Demonstrations}, 2021.

\end{thebibliography}
}

\clearpage
\onecolumn
\section*{Figures}
All figures are collected here after the 2-page narrative so the text pages remain uninterrupted and clearly within the page limit (figures excluded per the report instructions).

\begin{figure}[H]
\centering
\begin{subfigure}[t]{0.47\textwidth}
  \centering
  \maybeincludegraphics[width=\linewidth]{\figdir/01_performance_vs_method_sst2.png}
  \caption{SST-2}
\end{subfigure}\hfill
\begin{subfigure}[t]{0.47\textwidth}
  \centering
  \maybeincludegraphics[width=\linewidth]{\figdir/01_performance_vs_method_mrpc.png}
  \caption{MRPC}
\end{subfigure}
\caption{Performance comparison between full fine-tuning and LoRA ranks.}
\label{fig:perf-method}
\end{figure}

\begin{figure}[H]
\centering
\begin{subfigure}[t]{0.47\textwidth}
  \centering
  \maybeincludegraphics[width=\linewidth]{\figdir/02_performance_vs_trainable_params_sst2.png}
  \caption{SST-2}
\end{subfigure}\hfill
\begin{subfigure}[t]{0.47\textwidth}
  \centering
  \maybeincludegraphics[width=\linewidth]{\figdir/02_performance_vs_trainable_params_mrpc.png}
  \caption{MRPC}
\end{subfigure}
\caption{Performance versus trainable parameters. LoRA occupies the high-performance, low-parameter region.}
\label{fig:perf-params}
\end{figure}

\begin{figure}[H]
\centering
\begin{subfigure}[t]{0.47\textwidth}
  \centering
  \maybeincludegraphics[width=\linewidth]{\figdir/03_performance_vs_lora_rank_sst2.png}
  \caption{SST-2}
\end{subfigure}\hfill
\begin{subfigure}[t]{0.47\textwidth}
  \centering
  \maybeincludegraphics[width=\linewidth]{\figdir/03_performance_vs_lora_rank_mrpc.png}
  \caption{MRPC}
\end{subfigure}
\caption{Rank ablation. Classification performance improves at low rank and saturates around $r=8$.}
\label{fig:rank-ablation}
\end{figure}

\begin{figure}[H]
\centering
\begin{subfigure}[t]{0.47\textwidth}
  \centering
  \maybeincludegraphics[width=\linewidth]{\figdir/04_memory_savings_sst2.png}
  \caption{SST-2}
\end{subfigure}\hfill
\begin{subfigure}[t]{0.47\textwidth}
  \centering
  \maybeincludegraphics[width=\linewidth]{\figdir/04_memory_savings_mrpc.png}
  \caption{MRPC}
\end{subfigure}
\caption{Peak GPU memory comparison. LoRA reduces the memory required for adaptation.}
\label{fig:memory}
\end{figure}

\begin{figure}[H]
\centering
\begin{subfigure}[t]{0.47\textwidth}
  \centering
  \maybeincludegraphics[width=\linewidth]{\figdir/05_training_curves_sst2.png}
  \caption{SST-2}
\end{subfigure}\hfill
\begin{subfigure}[t]{0.47\textwidth}
  \centering
  \maybeincludegraphics[width=\linewidth]{\figdir/05_training_curves_mrpc.png}
  \caption{MRPC}
\end{subfigure}
\caption{Training curves for full fine-tuning and LoRA variants.}
\label{fig:curves}
\end{figure}

\begin{figure}[H]
\centering
\begin{subfigure}[t]{0.31\textwidth}
  \centering
  \maybeincludegraphics[width=\linewidth]{\figdir/paramsvsrank.png}
  \caption{Trainable params vs. rank}
\end{subfigure}\hfill
\begin{subfigure}[t]{0.31\textwidth}
  \centering
  \maybeincludegraphics[width=\linewidth]{\figdir/perfvslorarank.png}
  \caption{GLUE metrics vs. rank}
\end{subfigure}\hfill
\begin{subfigure}[t]{0.31\textwidth}
  \centering
  \maybeincludegraphics[width=\linewidth]{\figdir/perplexityvsrank.png}
  \caption{GPT-2 perplexity vs. rank}
\end{subfigure}
\caption{Cross-model summary of rank, parameter count, and language-modeling perplexity.}
\label{fig:summary}
\end{figure}

\end{document}